\documentclass[11pt,a4paper]{article}
\usepackage[utf8]{inputenc}
\usepackage[spanish]{babel}
\usepackage{hyperref}
\usepackage{geometry}
\usepackage{xcolor}
\geometry{margin=1in}

\begin{document}

\begin{titlepage}
    \centering
    \vspace*{0.5cm}
    {\fontsize{18}{22}\selectfont \textbf{UNIVERSIDAD TECNOLÓGICA DE QUERÉTARO} \par}
    \vspace{0.6cm}
    {\fontsize{14}{17}\selectfont División de Tecnologías de la Información y Comunicación \\ 
    Ingeniería en Desarrollo y Gestión de Software \par}
    \vspace{1.5cm}
    
    {\color{brown}\hrule height 2pt}
    \vspace{0.6cm}
    {\fontsize{20}{24}\selectfont \textbf{AEGIS Wiki - Elden Ring} \par}
    \vspace{0.4cm}
    {\fontsize{16}{20}\selectfont \textit{Documentación de Seguridad, Privacidad (LFPDPPP) y PWA (SA.4)} \par}
    \vspace{0.6cm}
    {\color{brown}\hrule height 2pt}
    
    \vspace{1.5cm}
    {\fontsize{12}{15}\selectfont
    \textbf{Asignatura:} Desarrollo para Dispositivos Inteligentes \\
    \vspace{0.2cm}
    \textbf{Evaluación:} Unidades 2 y 3: Integración del Ecosistema \\
    \vspace{0.2cm}
    \textbf{Nivel de Entrega:} SA - Satisfactorio \\
    \vspace{0.2cm}
    \textbf{Presenta:} Santiago Alberto Martínez Hernández \\
    \vspace{0.2cm}
    \textbf{Grupo / Cuatrimestre:} IDGS16 / Mayo - Agosto 2026 \par}
    
    \vfill
    {\large Santiago de Querétaro, Qro., 2026 \par}
\end{titlepage}

\newpage

\section{Introducción}
Esta documentación cubre los aspectos legales, de privacidad y la lista de comprobación de seguridad requeridos por la materia \textbf{Desarrollo para Dispositivos Inteligentes} para la \textbf{Evaluación 2 (Ecosistema Completo)}.

\section{Cumplimiento Legal y LFPDPPP}

\subsection{Identificación de la Ley Aplicable}
El ecosistema \textbf{AEGIS Wiki} cumple con la \textbf{Ley Federal de Protección de Datos Personales en Posesión de los Particulares (LFPDPPP)} en México. Dado que el backend de Express procesa e identifica cuentas de usuarios que residen en territorio nacional, el sistema establece controles claros sobre la privacidad y el uso ético de los datos.

\subsection{Datos Personales Recabados y Justificación}
El sistema procesa y almacena únicamente los siguientes datos personales:
\begin{enumerate}
    \item \textbf{Correo electrónico}: Obligatorio como identificador único de cuenta y para el canal de restablecimiento.
    \item \textbf{Contraseña}: Requerido para verificar la identidad en el login. Se procesa con cifrado unidireccional \texttt{bcrypt} en el servidor y nunca se almacena en texto plano.
    \item \textbf{Secreto TOTP (MFA)}: Clave alfanumérica compartida para el algoritmo de segundo factor (TOTP) generada por \texttt{otplib} para autenticación por proximidad o código manual.
\end{enumerate}

\textbf{Justificación}: La recolección de estos datos es mínima e indispensable para prestar el servicio de autenticación segura y personalización de las builds del usuario en el ecosistema.

\section{Aviso de Privacidad Básico (Simulado)}

\textbf{Responsable del Tratamiento}: El alumno del cuatrimestre Mayo-Agosto 2026 de la asignatura \textit{Desarrollo para Dispositivos Inteligentes}.\\

\textbf{Finalidad del Tratamiento}:\\
Los datos personales recabados serán utilizados exclusivamente para:
\begin{itemize}
    \item Evaluar el funcionamiento del inicio de sesión único (SSO) en múltiples pantallas.
    \item Establecer un canal seguro de comunicación (SSE/Bluetooth).
    \item Sincronizar las estadísticas del portador (ritmo cardíaco, horas de juego) entre el wearable, teléfono y Smart TV.
\end{itemize}

\textbf{Derechos ARCO (Acceso, Rectificación, Cancelación y Oposición)}:\\
El usuario puede solicitar en cualquier momento el acceso a sus datos de la bitácora, la rectificación de su nombre de perfil, la cancelación de su cuenta o la oposición al procesamiento. En esta plataforma académica, el usuario puede ejercer estos derechos enviando un correo al administrador o eliminando su cuenta de forma directa a través del módulo de autogestión.

\section{Plan de Retención y Eliminación de Datos}
\begin{itemize}
    \item \textbf{Servidor PostgreSQL}: Los registros de usuarios inactivos o suspendidos se marcan con eliminación lógica (\texttt{activo = false}) para preservar la integridad referencial de los históricos de auditoría, eliminándose físicamente a solicitud del titular de la cuenta.
    \item \textbf{Almacenamiento Local (PWA/Móvil)}: Los datos almacenados en el \texttt{localStorage} o \texttt{SharedPreferences} que excedan un período de \textbf{30 días} desde su creación son eliminados automáticamente al iniciar la aplicación.
    \item \textbf{Método de Eliminación}: El frontend de la Smart TV PWA evalúa en \texttt{ngOnInit} los timestamps de las claves con prefijo \texttt{aegis\_cache\_time\_} y elimina tanto el dato como el timestamp de expiración mediante la función \texttt{localStorage.removeItem()}.
\end{itemize}

\section{Checklist de Seguridad para la PWA}

\subsection{Content Security Policy (CSP)}
La PWA implementa restricciones estrictas en su cabecera \texttt{meta http-equiv="Content-Security-Policy"} en \texttt{index.html}:
\begin{itemize}
    \item \texttt{default-src 'self'}: Bloquea recursos externos no autorizados.
    \item \texttt{connect-src 'self' http://localhost:3000 https://aegis-wiki-backend.onrender.com}: Restringe el intercambio asíncrono y los Server-Sent Events (SSE) exclusivamente al backend del proyecto.
    \item \texttt{img-src} y \texttt{media-src}: Permite cargar imágenes y videos únicamente sobre canales cifrados \texttt{https:}.
\end{itemize}

\subsection{Validación de Origen (Event Origin)}
Para comunicaciones inter-pestañas mediante \texttt{BroadcastChannel}, se añade un middleware en Javascript que valida que \texttt{event.origin} coincida exactamente con el dominio base del proyecto (\texttt{https://wikieldenring.vercel.app}), descartando cualquier payload malicioso de inyección Cross-Site Scripting (XSS).

\subsection{Uso de HTTPS y Cifrado}
El backend en Render y la PWA en Vercel imponen de manera obligatoria la redirección de tráfico a canales seguros cifrados con el protocolo \textbf{HTTPS / TLS 1.3}. Se rechaza la transferencia sobre HTTP plano.

\subsection{Subresource Integrity (SRI)}
Se habilitó \texttt{"subresourceIntegrity": true} en \texttt{angular.json} para las compilaciones de producción de la PWA. Esto añade de forma automática firmas hashes criptográficas (\texttt{integrity="sha384-..."}) a todas las hojas de estilo y scripts cargados por el navegador, garantizando que los archivos de la app no hayan sido alterados o manipulados por terceros (ataques Man-in-the-Middle o inyecciones maliciosas) (\textbf{SA.4}).

\end{document}
